\documentclass[12pt]{article}
\usepackage[margin=1in]{geometry}
\usepackage{enumitem}
\usepackage{titlesec}
\usepackage{parskip}
\usepackage{hyperref}
\usepackage{fontenc}

\title{\textbf{Moore 1966 Claude Guide}\\
\large Moore, Barrington, Jr. \textit{Social Origins of Dictatorship and Democracy: Lord and Peasant in the Making of the Modern World}.\\
Boston: Beacon Press, 1966.}
\author{}
\date{}

\begin{document}

\maketitle
\tableofcontents
\newpage

%--------------------------------------------------
\section{Central Claims}
%--------------------------------------------------

\begin{itemize}[leftmargin=*, itemsep=4pt]
    \item There is no single path to the modern world. Societies arrive at modernity via distinct political-economic routes, and Moore identifies three: the \textbf{bourgeois-democratic route} (England, France, United States), the \textbf{capitalist-reactionary route} culminating in fascism (Germany, Japan), and the \textbf{communist route} (Russia, China).
    \item The decisive variable explaining which route a country takes is the character of the \textbf{commercialization of agriculture}---specifically, how the landed upper classes and the peasantry respond to the penetration of market relations into agrarian life, and what kind of coalition emerges between landed elites, the emerging bourgeoisie, and the peasantry.
    \item \textbf{``No bourgeois, no democracy''}: a economically and politically independent, sufficiently strong commercial/urban bourgeoisie---one not subordinated to or fused with a reactionary landed aristocracy---is a necessary condition for democratic development.
    \item The way the landed aristocracy commercializes agriculture matters enormously. If landlords commercialize by adapting to and working through market forces while retaining peasants on the land using repressive labor controls (labor-repressive agriculture), the result tends toward reactionary or fascist coalitions. If landlords are weakened, bought out, or transformed into capitalist farmers alongside a rising bourgeoisie, the result favors democracy.
    \item \textbf{Revolutionary violence, not the absence of conflict, underwrites gradualist and democratic outcomes.} England's ``gradualism'' was possible only because of prior revolutionary violence (the Civil War, enclosure) that broke the power of a reactionary peasantry/aristocracy alliance; there is no purely peaceful road to a modern democratic order.
    \item The strength or weakness of \textbf{peasant revolutionary potential} is itself variable and depends on the internal solidarity of the peasant community, the degree to which landlords are commercially or parasitically linked to the peasantry, and whether a revolutionary agent (party, army) can organize peasant grievances.
    \item \textbf{Fascism (``revolution from above'')} arises where a coalition of a weak, dependent bourgeoisie and a strong, labor-repressive landed elite pursues capitalist modernization without a genuine social revolution from below---commercializing the economy and strengthening the state while preserving reactionary social relations in the countryside.
    \item \textbf{Communist revolution} occurs where both the landed aristocracy and the bourgeoisie are too weak to modernize the agrarian order at all, leaving a large, exploited, and organizationally available peasantry as the social base for a revolutionary seizure of power by a disciplined party/military organization (this is elaborated further in Chapter 9 on the peasantry, though the China chapter---Chapter 4---remains central to the argument).
    \item Moore's method is explicitly \textbf{comparative-historical and structural}: he rejects both single-cause economic determinism and idealist/cultural explanations in favor of identifying structural conjunctures---especially class coalitions---that make certain political outcomes likely, without claiming strict determinism or universal laws.
    \item The book is a foundational text of \textbf{historical-comparative sociology} and directly inspired the "Moore thesis" research program (Skocpol, Paige, Luebbert, Downing, and others) on revolution, democratization, and regime transitions.
\end{itemize}

%--------------------------------------------------
\section{Core Case Studies}
%--------------------------------------------------

\begin{itemize}[leftmargin=*, itemsep=4pt]
    \item \textbf{England (Chapter 1)}: The paradigm case of the bourgeois-democratic route. Commercialization of agriculture (enclosure, wool trade) transformed the English aristocracy into a commercially-oriented class allied with, rather than opposed to, urban commercial interests; the Civil War and the ``violence'' of enclosure cleared the way for parliamentary government and gradual reform.
    \item \textbf{France (Chapter 2)}: A messier bourgeois-democratic route. Unlike England, the French peasantry survived the Revolution as small landholders, and the bourgeoisie's alliance with the peasantry against a parasitic (rather than commercially dynamic) aristocracy produced a more violent, discontinuous, and politically unstable path to democracy---revolution followed by reaction, empire, and repeated regime change.
    \item \textbf{China (Chapter 4)}: The gentry-landlord class was parasitic and bureaucratic rather than commercial, and failed to modernize agriculture; the imperial state decayed without a bourgeois or reformist alternative emerging. This vacuum, combined with an exploited and organizationally available peasantry, culminated in the Communist revolution led by the CCP, which organized peasant grievances the old order could not resolve.
    \item \textbf{Japan (Chapter 5)}: ``Revolution from above'' via the Meiji Restoration---a coalition of the imperial state bureaucracy and commercially-adapting landed elites modernized the economy and built a strong state while preserving repressive agrarian labor relations and excluding the peasantry from political power, culminating in fascism and militarism in the 1930s.
    \item \textbf{Germany (Chapter 8, treated alongside Japan)}: The Junker landed elite adapted to commercial agriculture using labor-repressive methods (the ``second serfdom'' east of the Elbe) and allied with a weak, dependent industrial bourgeoisie; the ``marriage of iron and rye'' produced a reactionary coalition that bypassed genuine bourgeois revolution and culminated in Nazism.
    \item \textbf{The United States (Chapter 3, not covered in this guide's chapter selection but relevant to Central Claims)}: The Civil War is analyzed as the ``last capitalist revolution''---the destruction of the labor-repressive plantation system in the South was a precondition for the unimpeded development of Northern industrial capitalism and democracy.
    \item \textbf{India (Chapter 6, not covered in this guide's chapter selection but relevant to Central Claims)}: A case of weak, incomplete modernization and democratization without social revolution---caste and village solidarity structures inhibited the emergence of either a strong peasant revolutionary movement or an aggressive commercializing landlord class, producing continued rural stagnation alongside formal parliamentary democracy.
\end{itemize}

%--------------------------------------------------
\section{Core Works Cited}
%--------------------------------------------------

\begin{itemize}[leftmargin=*, itemsep=4pt]
    \item \textbf{Karl Marx}---the foundational influence on Moore's class-analytic method; Moore's account of bourgeois revolution and class coalitions is directly indebted to (while also revising) Marxist historiography.
    \item \textbf{Max Weber}---on rationalization, bureaucracy, and the comparative sociology of world religions/civilizations; Moore's comparative-historical method and attention to the state bureaucracy (especially in the China and Japan chapters) is Weberian in spirit.
    \item \textbf{Alexis de Tocqueville}, \textit{The Old Regime and the French Revolution}---on the paradox that reform can precipitate revolution; informs Moore's treatment of the French case.
    \item \textbf{R.H. Tawney}, \textit{Religion and the Rise of Capitalism} and works on the English agrarian problem---background for the English enclosure and gentry commercialization story in Chapter 1.
    \item \textbf{Karl Polanyi}, \textit{The Great Transformation} (1944)---on the social dislocation caused by the commercialization of land and labor; a key intertext for Moore's emphasis on the disruptive effects of agrarian market penetration.
    \item \textbf{Eric Wolf}, \textit{Peasant Wars of the Twentieth Century} (1969)---a contemporary and complementary work on peasant revolution, often read alongside Moore.
    \item \textbf{Theda Skocpol}, \textit{States and Social Revolutions} (1979)---the most influential critical successor to Moore; shifts emphasis from class coalitions to state structures and international/military pressures, explicitly building on and revising the ``Moore thesis.''
    \item \textbf{Jeffery Paige}, \textit{Agrarian Revolution} (1975)---extends Moore's agrarian class-coalition framework to twentieth-century Third World revolutions using a more systematic comparative-quantitative method.
    \item \textbf{Gregory Luebbert}, \textit{Liberalism, Fascism, or Social Democracy} (1991)---directly engages and revises Moore's account of interwar European regime divergence, emphasizing the role of the peasantry and working-class/agrarian coalitions.
    \item \textbf{Brian Downing}, \textit{The Military Revolution and Political Change} (1992)---revises Moore's account by emphasizing the impact of early modern military competition and state-building on subsequent regime trajectories.
    \item \textbf{Alexander Gerschenkron}, \textit{Economic Backwardness in Historical Perspective} (1962)---on late industrialization and the role of the state in substituting for a weak bourgeoisie; relevant background for the Germany and Japan ``revolution from above'' cases.
\end{itemize}

%--------------------------------------------------
\section{Chapter 1: England and the Contributions of Violence to Gradualism}
%--------------------------------------------------

\begin{itemize}[leftmargin=*, itemsep=4pt]
    \item \textbf{The puzzle of English gradualism}: England is usually cited as the paradigm of peaceful, evolutionary political development. Moore argues this reputation obscures the substantial violence---enclosure, civil war, agrarian dispossession---that made gradualism possible; ``gradualism'' describes the nineteenth and twentieth centuries only because the seventeenth and eighteenth centuries were not gradual at all.
    \item \textbf{Commercialization of agriculture via enclosure}: English landlords responded to growing wool and grain markets by enclosing common land and converting the peasantry's customary rights into commercial tenancies, transforming the aristocracy into a commercially oriented class with interests aligned with, not opposed to, urban commerce.
    \item \textbf{The Civil War (1640s) as bourgeois revolution}: Moore reads the English Civil War as a crucial---if incomplete and socially conservative---bourgeois revolution: it broke the political power of the crown and cleared obstacles (feudal dues, royal prerogative over commerce) to the development of agrarian and commercial capitalism, even though it was fought in religious and constitutional language.
    \item \textbf{Destruction of the peasantry as a class}: Unlike France, the English peasantry was effectively eliminated as an independent political and economic force well before the age of democratic reform (through enclosure and the growth of agrarian wage labor)---removing peasants as a potential base for either reactionary or revolutionary mobilization in the nineteenth century.
    \item \textbf{Repeal of the Corn Laws (1846)}: The decisive triumph of urban-industrial and commercial-agrarian interests over the last vestige of protectionist landed power, symbolizing the completed alliance between commercial landlords and the industrial bourgeoisie.
    \item \textbf{Labor-intensive vs.\ labor-repressive agriculture}: English commercial agriculture relied on market discipline and wage labor (eviction, enclosure) rather than repressive, coercive labor controls (serfdom, forced labor)---a crucial contrast with the German and Eastern European path.
    \item \textbf{Parliament as the arena of elite accommodation}: Because the English aristocracy commercialized successfully and the peasantry had been eliminated as an independent class, later extensions of the franchise (1832, 1867, 1884) could proceed gradually through parliamentary bargaining rather than revolutionary rupture.
    \item \textbf{General lesson}: Gradualist, reformist democratic development in the industrial era is not the absence of violent conflict but its historical prerequisite---violence earlier in the process (breaking reactionary agrarian power, eliminating the peasantry as an independent class) is what permits gradualism later.
\end{itemize}

%--------------------------------------------------
\section{Chapter 2: Evolution and Revolution in France}
%--------------------------------------------------

\begin{itemize}[leftmargin=*, itemsep=4pt]
    \item \textbf{A messier bourgeois route}: France reaches democracy via the same broad ``bourgeois'' route as England, but far more violently and discontinuously, because the underlying agrarian class structure differed in crucial respects.
    \item \textbf{The French aristocracy as parasitic, not commercial}: Unlike the English gentry, the French nobility largely failed to commercialize agriculture directly; instead, it relied on royal privilege, tax exemptions, and seigneurial dues extracted from a peasantry that retained land rights---producing an aristocracy dependent on the crown rather than the market.
    \item \textbf{Survival of the peasantry as smallholders}: Because the French peasantry was not dispossessed the way the English peasantry was, it survived the Revolution largely as a class of small landowners. The Revolution destroyed seigneurial dues and redistributed some land but did not create large-scale commercial agriculture; France remained a country of small peasant proprietors well into the twentieth century.
    \item \textbf{The Revolution (1789) as bourgeois revolution}: The Revolution smashed the fused aristocratic-monarchical order and cleared legal and institutional obstacles to capitalism (guilds, internal tariffs, seigneurial privilege), but its social base---an alliance of the bourgeoisie with a smallholding peasantry---produced a more conservative economic outcome (a nation of small farms and small enterprises) than England's more thoroughgoing agrarian capitalism.
    \item \textbf{Peasant conservatism after the Revolution}: Having won secure property rights in 1789, the French peasantry became a fundamentally conservative political force in the nineteenth century, wary of further social upheaval that might threaten smallholding---famously providing the rural base for both Napoleonic plebiscitary rule and later right-wing and Bonapartist politics.
    \item \textbf{Recurrent instability}: The absence of a stable elite coalition (unlike England's fused landlord-bourgeois alliance) produced a recurring cycle of revolution, reaction, and regime change (1789, 1830, 1848, 1851, 1871) rather than steady, gradual reform.
    \item \textbf{Comparison with England}: Where England achieved bourgeois dominance through an alliance of commercializing landlords and urban capital (with the peasantry eliminated), France achieved it through revolutionary destruction of a parasitic aristocracy while preserving the peasantry---two structurally different roads converging, eventually, on democratic outcomes.
\end{itemize}

%--------------------------------------------------
\section{Chapter 4: The Decay of Imperial China and the Origins of the Communist Variant}
%--------------------------------------------------

\begin{itemize}[leftmargin=*, itemsep=4pt]
    \item \textbf{The gentry as bureaucratic-parasitic, not commercial}: Chinese landlords (the gentry) derived their status primarily from ties to the imperial bureaucracy and examination system, not from entrepreneurial commercialization of agriculture; they extracted rent and taxes but did not transform agricultural production the way English or even Prussian landlords did.
    \item \textbf{Absence of an autonomous bourgeoisie}: China lacked a commercial/urban class capable of contesting gentry-bureaucratic dominance; merchants were subordinated to and dependent on the imperial state and the gentry rather than constituting an independent political force.
    \item \textbf{Failure of agrarian modernization}: Because the gentry had no interest in capitalist transformation of agriculture (their income depended on traditional rent extraction and office-holding, not market competitiveness), Chinese agriculture stagnated technologically and organizationally even as population pressure intensified.
    \item \textbf{Imperial state decay}: The nineteenth-century collapse of the imperial order (Taiping Rebellion, foreign imperialism, the fall of the Qing in 1911) was not accompanied by the emergence of either a reformist landlord class or a strong bourgeoisie capable of building a modern capitalist state, as had happened (in different ways) in England, France, and Japan.
    \item \textbf{The warlord/Guomindang interregnum}: The Nationalist (Guomindang) government is portrayed as unable or unwilling to break gentry landlordism and build a genuine reformist coalition---it remained tied to urban and landlord interests and failed to address the peasant question, leaving revolutionary space open.
    \item \textbf{Peasant exploitation and organizational availability}: A large, exploited, and increasingly desperate peasantry, no longer buffered by traditional village solidarities that were breaking down under commercial and demographic pressure, became available for revolutionary mobilization by a disciplined outside organization.
    \item \textbf{The Chinese Communist Party as organizing agent}: The CCP succeeded where the Guomindang failed by directly organizing peasant grievances (land redistribution, rent reduction) into a coherent revolutionary movement, providing the leadership and discipline the peasantry could not generate spontaneously---prefiguring the argument developed further in Chapter 9.
    \item \textbf{General lesson}: Communist revolution in China resulted from the absence of \emph{any} class capable of modernizing agriculture and the state from above (neither a Junker-style commercializing gentry nor an autonomous bourgeoisie), leaving revolution from below as the only route to a modern state.
\end{itemize}

%--------------------------------------------------
\section{Chapter 5: Asian Fascism: Japan}
%--------------------------------------------------

\begin{itemize}[leftmargin=*, itemsep=4pt]
    \item \textbf{Revolution from above via the Meiji Restoration (1868)}: Confronted with the threat of Western imperialism, a coalition of lower samurai, court nobles, and elements of the commercial class overthrew the Tokugawa shogunate and built a modern, centralized state without a genuine social revolution against the landed elite.
    \item \textbf{Landlord-state alliance}: Unlike England, Japanese landlords retained and even strengthened their position over the peasantry during modernization; the Meiji state relied on land taxes extracted from a subordinated peasantry to finance industrialization, preserving repressive agrarian social relations while building a modern economy and military.
    \item \textbf{Labor-repressive agriculture}: Japanese landlords used tenancy arrangements, debt bondage, and high rents (extracted in kind) to control the peasantry, extracting the agricultural surplus needed for state-led industrialization without transforming property relations in a capitalist direction.
    \item \textbf{Weak, dependent bourgeoisie}: Japanese industrial capitalists (the \textit{zaibatsu}) developed under heavy state sponsorship and protection rather than as an independent political force contesting landlord or bureaucratic power; they became junior partners in, rather than challengers to, the landlord-bureaucratic-military coalition.
    \item \textbf{Exclusion of the peasantry from political power}: Peasants were mobilized economically (as a source of surplus and later as army recruits) but excluded from meaningful political participation, unlike the eventual (if delayed) political incorporation of the working and agrarian classes in the Western democracies.
    \item \textbf{The road to fascism/militarism in the 1930s}: The absence of a bourgeois revolution against landlordism, combined with the weakness of parliamentary institutions and the political prominence of the military and landed interests, culminated in the militarist, ultranationalist regime of the 1930s and 1940s---Moore's ``fascism from above.''
    \item \textbf{Comparison with Germany}: Japan and Germany both industrialize via a coalition of a traditional landed elite and a weak bourgeoisie under strong state sponsorship (``revolution from above''), converging on reactionary/fascist outcomes despite very different cultural and historical starting points---evidence for Moore's structural, not culturalist, mode of explanation.
\end{itemize}

%--------------------------------------------------
\section{Chapter 7: The Democratic Route to Modern Society}
%--------------------------------------------------

\begin{itemize}[leftmargin=*, itemsep=4pt]
    \item \textbf{Synthesizing the bourgeois-democratic cases}: This chapter draws together the England, France, and (briefly) United States cases to specify the general structural conditions under which capitalist modernization produces parliamentary democracy rather than reaction or revolution.
    \item \textbf{``No bourgeois, no democracy''}: The single most quoted formulation in the book---a economically independent and politically assertive bourgeoisie, not subordinated to a reactionary landed class, is indispensable to democratic outcomes.
    \item \textbf{Four preconditions for the democratic route}: (1) development of a balance of power between the crown/central authority and the landed aristocracy that prevents excessive royal absolutism; (2) a turn toward commercial agriculture that weakens rather than reinforces repressive labor controls; (3) weakening or destruction of the landed aristocracy as an independent political class (or its fusion with the bourgeoisie); (4) prevention of an aristocratic-peasant coalition against the bourgeoisie/commercial classes.
    \item \textbf{Elimination of a revolutionary peasantry}: In the successful democratic cases, either the peasantry is eliminated as an independent class (England) or is transformed into a conservative smallholding class with a stake in the existing property order (France, and the American North's yeoman farmers)---removing peasants as a base for either reactionary or revolutionary mobilization.
    \item \textbf{The bourgeois revolution need not be fully ``bourgeois'' in ideology}: Moore stresses that the Puritan Revolution and the French Revolution were not fought explicitly for capitalism, yet both had the structural effect of clearing away the political-legal remnants of the agrarian-aristocratic order, which is what matters causally.
    \item \textbf{Democracy requires the destruction, not merely the reform, of the old agrarian order}: Piecemeal reform that preserves the political power of a labor-repressive landed class does not produce democracy; some decisive rupture (civil war, revolution) with the old order's political power is necessary, even if economically conservative.
    \item \textbf{Contingency and multiple causation}: Moore explicitly resists a single-variable determinism; the ``route'' is a structural tendency produced by a specific configuration of class forces and historical timing, not an iron law---different sequencing or external shocks could, in principle, produce different outcomes.
\end{itemize}

%--------------------------------------------------
\section{Chapter 8: Revolution from Above and Fascism}
%--------------------------------------------------

\begin{itemize}[leftmargin=*, itemsep=4pt]
    \item \textbf{Defining ``revolution from above''}: A pattern of capitalist modernization carried out by a coalition of the state bureaucracy, the military, and a landed elite, without a genuine social revolution against agrarian reaction---contrasted with the ``revolution from below'' path of England and France.
    \item \textbf{Germany as the paradigm European case}: The Junker landlords of Prussia commercialized agriculture east of the Elbe using labor-repressive methods (the ``second serfdom''), while allying with a weak, state-dependent industrial bourgeoisie---the famous ``marriage of iron and rye''---to modernize the economy while preserving reactionary political and agrarian structures.
    \item \textbf{Bismarckian unification as elite-managed modernization}: German unification (1871) and subsequent industrialization were achieved through state-led, top-down policy (tariffs, social insurance, military buildup) that co-opted rather than confronted the landed aristocracy, foreclosing the possibility of a bourgeois-democratic revolution.
    \item \textbf{Weakness of German liberalism}: The German bourgeoisie, dependent on state contracts, protective tariffs, and alliance with the Junkers for political stability, never mounted an independent challenge to aristocratic political power comparable to the English or French bourgeoisie---the failure of the 1848 revolutions is treated as paradigmatic of this weakness.
    \item \textbf{From authoritarian conservatism to fascism}: The unresolved tension between an increasingly industrial society and a political order still dominated by pre-industrial landed and military elites created the structural opening for Nazism, which mobilized middle-class and rural resentment while ultimately serving the interests of the traditional landed and industrial elites.
    \item \textbf{Comparison with Japan}: Moore treats Germany and Japan as structurally parallel instances of ``revolution from above''---despite vast cultural differences, both combine a labor-repressive landed elite, a weak/dependent bourgeoisie, and an assertive state/military apparatus, converging on fascist outcomes.
    \item \textbf{Fascism as a modernizing, not merely reactionary, phenomenon}: Fascist regimes are not simple throwbacks to the pre-industrial past; they represent a distinct route to industrial modernity that preserves reactionary social relations in agriculture while aggressively modernizing industry and the military---a ``capitalist reactionary'' rather than democratic-capitalist path.
\end{itemize}

%--------------------------------------------------
\section{Chapter 9: The Peasants and Revolution}
%--------------------------------------------------

\begin{itemize}[leftmargin=*, itemsep=4pt]
    \item \textbf{The central question}: Under what conditions does the peasantry become a revolutionary---rather than conservative or politically passive---force? Moore treats the peasant question as analytically distinct from, but essential to, the class-coalition arguments of the earlier chapters.
    \item \textbf{Peasant solidarity as a precondition for revolt}: Peasant communities with strong internal cohesion (village solidarity, shared land tenure institutions, mutual economic dependence) are more capable of collective action against landlords or the state than atomized or hierarchically stratified peasant societies.
    \item \textbf{The role of the landlord's relationship to the peasantry}: Where landlords are physically and socially present in the village and personally responsible for extracting the surplus (as in parts of China and Russia), peasant hostility can be focused directly on a visible local target, facilitating revolt. Where landlords are absentee or where extraction is mediated by commercial/market mechanisms, peasant grievances are more diffuse and harder to mobilize.
    \item \textbf{Commercialization can either dampen or inflame peasant revolt}: Market penetration that preserves communal solidarity while increasing exploitation (e.g., commercialized rent extraction without dispossession) can heighten peasant grievance and organizational capacity; market penetration that dissolves the peasant community altogether (e.g., English-style enclosure) can eliminate the peasantry as a political actor entirely.
    \item \textbf{The need for outside leadership/organization}: Even highly aggrieved and solidaristic peasantries rarely produce sustained, large-scale revolution without organization and leadership from an external (usually urban-based) political or military organization capable of providing strategy, discipline, and a program (land reform)---as with the Bolsheviks in Russia and the CCP in China.
    \item \textbf{Comparative cases revisited}: Russia and China are treated as the paradigm cases of peasant-based revolution, where a weak or absent bourgeoisie and a parasitic, non-commercializing landlord class left the peasantry as the only available social base for revolutionary transformation, given effective outside leadership.
    \item \textbf{Peasant revolution as a route to modernity, not backwardness}: Like ``revolution from above,'' peasant revolution is presented as a distinct, historically important route to a modern (industrial, bureaucratic) social order---not a simple failure to modernize, but modernization achieved through a different class alliance and a different, more violent, mechanism.
    \item \textbf{Synthesis with the book's overall argument}: This chapter completes Moore's three-route typology by specifying the peasant-organizational conditions underlying the communist route (China, Chapter 4; Russia, discussed comparatively throughout), paralleling the bourgeois conditions specified for the democratic route (Chapter 7) and the landlord-bureaucratic conditions specified for the fascist route (Chapter 8).
\end{itemize}

\end{document}
